\section{Navigation for Source Section III}
\label{sec:numerical-results-overview}

\sourceeqrange{Eqs. (3.1)--(3.20) and source figures 1--15}

\subsection*{Purpose}

Source Section III is a simulation section.  The companion treats it as a
reading guide: it identifies the nondimensional variables, boundary branches,
scenario assumptions, and scalar estimates needed to read the published figures.
It does not reconstruct the simulation code, regenerate the figures, or infer
which unpublished simulation-code branch was used for the Appendix-B pressure
tensor.  It does not validate the numerical results.
Here ``scenario'' means a source-described numerical example, not a reproduced
run in this repository.

\subsection*{Common Method Branch}

All scenarios in this section inherit the Section III.A specialization
\(K=0\) with constant \(C\), the length scale \(\ell\) in Eq. (3.1), the wall
density condition in Eq. (3.2), density-proportional transport coefficients in
Eqs. (3.3)--(3.4), and the time scale \(t_0=\ell^2/\nu_0\) in Eq. (3.5).  The
velocity is no-slip on the computational walls, so the mean number density is a
conserved simulation input.  Thermal boundary data vary by scenario: some runs
use one wall temperature on all boundaries, while heat-flow runs prescribe
bottom and top temperatures and insulating side walls.


\subsection*{How to Read a Scenario}

The companion uses the same four-step reading pattern for every source
scenario.  First identify the method branch inherited from Section III.A:
scaled variables, transport coefficients, no-slip velocity, density-wall data,
and thermal boundary data.  Second identify which local balance law from
Section II.D is being used as background.  Third read any scalar estimate, such
as \(\sigma\), \(g^*\), \(v_c\), \(\lambda_{\rm eff}\), or \(Nu\), as a typed
scale or diagnostic.  Fourth treat the published figures as source simulation
outputs.  The fourth step is intentionally not reproduced here.

\subsection*{Scenario Map}

\begin{description}[leftmargin=0.19\linewidth,style=nextline]
\item[III.A, method]
Eqs. (3.1)--(3.11) and Appendix B provide scales, wall conditions, and estimate
variables.
\item[III.B, adiabatic expansion]
Eq. (3.12) and figures 1--2 are read as a boundary-cooling scenario; no time
trace is reproduced.
\item[III.C, piston effect]
Eqs. (3.13)--(3.16) and figures 3--4 organize pressure, temperature, and
latent-heat estimates.
\item[III.D, two-phase heat flow]
Eq. (3.16) and figures 5--7 connect droplet heat-flow interpretation to the
latent-heat velocity scale.
\item[III.E, steady heat conduction]
Eqs. (3.17)--(3.19) and figures 8--10 define effective conductivity and a
Nusselt-style ratio.
\item[III.F, boiling in gravity]
Eq. (3.20) and figures 11--13 record the gravity branch and the adiabatic
bulk-temperature-gradient estimate.
\item[III.G, wetting dynamics]
Eq. (3.2), Eq. (3.18), and figures 14--15 connect wetting dynamics to the wall
condition and bottom heat-flux diagnostic.
\end{description}

\begin{figure}[htbp]
\centering
\resizebox{0.92\linewidth}{!}{\begin{tikzpicture}[
  font=\small,
  node distance=7mm and 7mm,
  box/.style={draw=black, line width=0.45pt, rounded corners=2pt,
    align=center, text width=31mm, minimum height=10mm, fill=gray!8},
  est/.style={draw=black, line width=0.45pt, rounded corners=2pt,
    align=center, text width=34mm, minimum height=10mm, fill=gray!14},
  open/.style={draw=black, dashed, line width=0.45pt, rounded corners=2pt,
    align=center, text width=36mm, minimum height=10mm, fill=white},
  arrow/.style={-{Latex[length=2.2mm]}, line width=0.5pt},
  darrow/.style={-{Latex[length=2.2mm]}, dashed, line width=0.5pt}
]
\node[est] (method) {III.A method\\scales and boundaries};
\node[box, below left=of method] (exp) {III.B\\adiabatic expansion};
\node[box, below=of method] (piston) {III.C--D\\piston and heat flow};
\node[box, below right=of method] (steady) {III.E--G\\steady heat, boiling, wetting};
\node[est, below=of piston] (checks) {Finite checks\\units, ratios, and scales};
\node[open, below=of checks] (no) {No figure reproduction\\no code inference};

\draw[arrow] (method) -- (exp);
\draw[arrow] (method) -- (piston);
\draw[arrow] (method) -- (steady);
\draw[arrow] (exp) -- (checks);
\draw[arrow] (piston) -- (checks);
\draw[arrow] (steady) -- (checks);
\draw[darrow] (checks) -- (no);
\end{tikzpicture}
}
\caption{Section III scenario-navigation map.  The companion checks finite
scales and units used to read the source scenarios; it does not reproduce
figures, regenerate time traces, or infer unpublished implementation details.}
\label{fig:section-iii-scenario-navigation}
\end{figure}

\subsection*{What This Companion Verifies}

The verification layer checks only finite formulas: equivalence of the two
forms of \(\sigma\), the density-proportional transport ratio, dimensionlessness
of \(g^*\), consistency of the Eq. (3.9) sound-scale denominator, the acoustic
traversal scale, the capillary-wave estimate, the inverse-compressibility
identity, the piston-effect scalar estimates, units of the latent-heat velocity
estimate, units of the effective conductivity definition, the one-phase Nusselt
estimate, the gravity-gradient order scale, the heat-flux normalization, and
the scaled wall-density condition.  These checks do not validate the numerical
results, interface dynamics, boiling transition, or wetting kinetics.
They are companion estimate checks attached to source equation labels, not
simulation benchmarks.

The executable mirrors for the Section III finite-estimate checks are:
\begin{center}
\small
\begin{tabular}{@{}l@{}}
\texttt{scripts/python/section\_iii\_simulation\_estimates.py}\\
\texttt{scripts/wolfram/section\_iii\_simulation\_estimates\_checks.wls}\\
\texttt{tests/test\_section\_iii\_simulation\_guide.py}
\end{tabular}
\end{center}
Those files check scalar units, ratios, and scale identities only; they do not
rerun or validate the source simulations.

\subsection*{Source-Procedure Completion Status}

The Section III rows assigned to this completion task are covered as
source-reading derivations, not as numerical reruns:
\begin{itemize}
\item Section III.A supplies the shared \(K=0\), constant-\(C\) branch,
  density-wall condition, transport closures, and finite scales
  \(\ell,t_0,\sigma,g^*,c,A_s,\omega_q\).
\item Sections III.B--III.G record typed scalar estimates and diagnostics:
  inverse compressibility, piston-effect pressure and temperature estimates,
  latent-heat velocity, effective conductivity, bottom heat flux, Nusselt
  ratio, and gravity/adiabatic-gradient estimates.
\item Published figures, time traces, boiling transitions, droplet paths, and
  wetting dynamics remain source outputs.  The companion does not reproduce or
  validate them.
\end{itemize}
The unresolved Appendix-B simulation-code branch and the global wall-transfer
map remain outside these finite estimate checks.

\subsection*{Tagged verification complement}

\OnukiLinkSectionIII{The tagged Section III complement} collects the reading
route, executable scalar checks, and negative controls. It does not provide a
numerical reproduction or validation of the source simulations.

\subsection*{Open Issues Carried Forward}

The unpublished simulation code remains unavailable.  The dimensional-source
and printed-B3 pressure-tensor branches therefore remain separate from the
simulation-code branch.  Global wall and surface-energy transfers are still
model-specific boundary data, not consequences of the scalar checks below.
